\section{Overall progress}\label{sec:scope}

The investigation has one common objective: decide whether positive total logarithmic Mahler measures of integer polynomials are separated from zero.
The four strategies below address that objective through different structures.
Their intermediate achievements must be compared with the same unrestricted quantifiers.

\begin{definition}\label{sec:formulations}
For a nonzero integer polynomial $P(t)=a_d\prod_{i=1}^d(t-\alpha_i)$, its Mahler measure and logarithmic Mahler measure are
\begin{equation}\label{fo:mahler}
 M(P)=|a_d|\prod_{i=1}^d\max\{1,|\alpha_i|\},\qquad m(P)=\log M(P).
\end{equation}
Roots are counted with multiplicity, and logarithms are natural.
\end{definition}
We use the usual convention of \citet[\S1]{Smy08}.
The target is to determine whether
\begin{equation}\label{fo:gap}
 \mathfrak L=\inf\{m(P):P\in\Z[t]\setminus\{0\},\ m(P)>0\}
\end{equation}
is positive or zero, with degree and coefficient size unrestricted.
Here $m(P)$ is total measure, not the degree-normalized height $m(P)/d$.
The investigation remains mathematically unresolved and is presently paused after the completed finite norm-one comparison.

The strongest general arithmetic change is that positive auxiliary surplus is available on every nontorsion algebraic unit.
To state it, for $F(t)=\sum_jf_jt^j\ne0$ put
\[
 S_2(F)=\Bigl(\sum_jf_j^2\Bigr)^{1/2},\qquad
 S_1(F)=\sum_j|f_j|,\qquad \rho(F)=S_1(F)^2/S_2(F)^2.
\]
These are norms of the complete coefficient vector.
A nontorsion algebraic unit is an algebraic integer with integral inverse which is not a root of unity.

\begin{theorem}\label{ar:uniform-auxiliary}
Let $P\in\Z[t]$ be monic and irreducible of degree $d$, with constant term $\pm1$ and no root of unity, and let $\alpha_1,\ldots,\alpha_d$ be all its roots.
There are a primitive nonconstant $F\in\Z[t]$, integers $r\ge1$ and $0\le j<2d$ such that
\begin{align}
 \deg F&\le2d-1,&
 F(\alpha_i)&=2^r\alpha_i^j\quad(1\le i\le d),\notag\\
 |\Res(P,F)|&=2^{rd},&
 \log(2^r)-\log S_2(F)&\ge\log(2/\sqrt3),\qquad \rho(F)\le2d.
 \label{ar:dyadic-output}
\end{align}
Allowing an arbitrary positive integer $Q$ instead of a power of two, the surplus can be at least $\frac12\log2$.
\end{theorem}

The proof in Section~\ref{ar:auxiliary} uses a finite rational projection, integer rounding and primitive normalization.
This theorem separates existence of positive surplus from the remaining cost of using that surplus at all conjugates.
Signed aggregation gives a square-root degree bound on the latter cost, but the unrestricted problem requires a useful input-uniform estimate.
The exact finite comparisons show that genuine unit directions can improve optimized monomial baselines; they do not establish varying-input coverage.

Table~\ref{ov:routes} summarizes the four strategies.
The subsequent sections explain their results, proof mechanisms and unresolved interfaces.
Two subsidiary strategies sit inside algebraic entropy, one inside linear-group growth, and one inside arithmetic.

\begin{table}[htbp]
\centering\small
\begin{tabular}{@{}>{\raggedright\arraybackslash}p{.18\textwidth}
 >{\raggedright\arraybackslash}p{.37\textwidth}
 >{\raggedright\arraybackslash}p{.37\textwidth}@{}}
\toprule
Strategy & Principal established progress & Remaining interface\\
\midrule
Algebraic entropy
& Rational primary decomposition; the quantitative Shannon comparison
and conditional critical Fourier regularity; complete costs for two
integer completion families.
& A uniform orbit constraint or a finite-capture mechanism that survives
unbounded complexity.\\[5pt]
Linear-group growth
& The original Lehmer return core is infinite; stationary laws have a
necessary diffusion scale; paid block certificates and Gaussian
comparison losses have precise fixed-input bounds.
& Control of the complete original weighted law, especially neutral
coordinates, uniformly over varying small-measure inputs.\\[5pt]
Group determinants
& An exact scalar return for square Laurent weak isomorphisms; full
companion and coupling costs are retained in integer constructions.
& A useful completion or operator estimate with every extra scalar
factor and its full cost controlled.\\[5pt]
Arithmetic
& Uniform positive auxiliary surplus, a signed bundle inequality and
square-root projection cost; five exact finite optima and two strict
unit-only improvements.
& Uniform aggregate correlation control, or same-exponent linear
cyclic-norm income, on all sufficiently small-measure inputs.\\
\bottomrule
\end{tabular}
\caption{Progress and the missing uniformity in each strategy.
Conditional statements and finite certificates retain their stated hypotheses and domains.}
\label{ov:routes}
\end{table}

\paragraph{What the current results decide.}\label{sec:outlook}
The original degree-ten return-core theorem concerns one fixed input.
The block-certificate bounds concern a specified way of obtaining a lower covariance bound.
The completion theorems cover specified complete families of integer matrices.
The five finite optima concern explicitly retained direction lists; the sixth enlarged optimum remains uncertified.
These distinct scopes do not prove a positive value of $\mathfrak L$ or construct positive measures tending to zero.

\paragraph{Proof and review scope.}
The report gives complete elementary reductions and the principal arithmetic arguments, and explicitly identifies outlines of longer adopted proofs.
Results have recorded reviews of their actual inputs, with self-review and separately organized AI review distinguished.
This report does not assert independent human certification, formal verification or novelty for every intermediate argument.
Mathematical sources are cited where used.
The standalone exact-arithmetic supplement verifies the finite comparison, including the retained failed dual; its success does not certify an optimum beyond those finite families.
